\section{Теорема о пределе сложной функции. Первый замечательный предел. (Beta)}

\subsection{Теорема о пределе сложной функции}

\begin{theorem}
Пусть функция $g(x)$ определена в проколотой окрестности точки $x_0$ и
\[
\Lim{x}{x_0} g(x)=a.
\]
Пусть функция $f(t)$ определена в проколотой окрестности точки $a$ и
\[
\Lim{t}{a} f(t)=A.
\]
Тогда, если $g(x)\neq a$ при $x\neq x_0$ (достаточно близких к $x_0$),
выполняется
\[
\Lim{x}{x_0} f(g(x))=A.
\]
\end{theorem}

\begin{Proof}
Пусть $x_n \to x_0$, $x_n \neq x_0$.
Тогда $g(x_n)\to a$.
Так как $\Lim{t}{a} f(t)=A$, то
\[
f(g(x_n)) \to A.
\]
По определению предела функции по Гейне получаем
\[
\Lim{x}{x_0} f(g(x))=A.
\]
\end{Proof}

\begin{remark}
Если функция $f$ непрерывна в точке $a$, то условие можно записать проще:
если
\[
\Lim{x}{x_0} g(x)=a,
\]
то
\[
\Lim{x}{x_0} f(g(x))=f(a).
\]
\end{remark}

\subsection{Первый замечательный предел}

\begin{theorem}[Первый замечательный предел]
\[
\Lim{x}{0}\frac{\sin x}{x}=1.
\]
\end{theorem}

\begin{Proof}
Для $x>0$ (в радианах) справедливо неравенство:
\[
\sin x < x < \tan x.
\]
Разделим все части на $\sin x$:
\[
1 < \frac{x}{\sin x} < \frac{1}{\cos x}.
\]
Переходя к обратным величинам:
\[
\cos x < \frac{\sin x}{x} < 1.
\]
При $x \to 0$ имеем $\cos x \to 1$.
По теореме о зажатой функции:
\[
\frac{\sin x}{x} \to 1.
\]
Аналогично для $x<0$.
\end{Proof}

\subsection{Следствия первого замечательного предела}

\begin{enumerate}
\item
\[
\Lim{x}{0}\frac{\sin(kx)}{kx}=1,
\qquad
\Lim{x}{0}\frac{\sin(kx)}{x}=k.
\]

\item
\[
\Lim{x}{0}\frac{\tan x}{x}=1.
\]

\item
\[
\Lim{x}{0}\frac{1-\cos x}{x}=0.
\]

\item
\[
\Lim{x}{0}\frac{1-\cos x}{x^2}=\frac12.
\]
\end{enumerate}

\subsection{Примеры}

\begin{example}
Найти предел:
\[
\Lim{x}{0}\frac{\sin 3x}{x}.
\]

\[
\frac{\sin 3x}{x}
=
\frac{\sin 3x}{3x}\cdot 3
\to 1\cdot 3=3.
\]
\end{example}

\begin{example}
Найти предел:
\[
\Lim{x}{0}\frac{1-\cos x}{x^2}.
\]

Используем тождество:
\[
1-\cos x = 2\sin^2\frac{x}{2}.
\]

Тогда
\[
\frac{1-\cos x}{x^2}
=
\frac{2\sin^2\frac{x}{2}}{x^2}
=
2\left(\frac{\sin\frac{x}{2}}{\frac{x}{2}}\right)^2
\cdot
\frac{1}{4}
\to
\frac12.
\]
\end{example}
